\section{Knowledge Validation and Code Verification}
The non-deterministic nature of Large Language Models introduces significant risks when they are used to generate critical software artifacts, such as microcontroller firmware or optimization scripts. To ensure the reliability of an agentic framework, outputs must be validated through systematic evaluation frameworks and automated verification loops.

\subsection{DeepEval: Quantitative RAG Evaluation}
For systems relying on Retrieval-Augmented Generation, measuring the quality of the retrieved context and the generated response is essential. \textit{DeepEval} is an open-source framework that implements the "LLM-as-a-Judge" paradigm to provide quantitative metrics for RAG performance \cite{DeepEval2024}. Unlike traditional metrics (e.g., BLEU or ROUGE), DeepEval evaluates semantic alignment across four primary dimensions:
\begin{itemize}
    \item \textbf{Faithfulness:} Measures whether the generated claims are factually supported by the retrieved context, preventing hallucinations.
    \item \textbf{Answer Relevancy:} Evaluates how well the response addresses the original user query.
    \item \textbf{Contextual Relevancy:} Assesses the quality of the retrieval process by measuring the actual utility of the fetched documents.
    \item \textbf{Hallucination Detection:} Specifically flags claims that contradict the provided source material.
\end{itemize}

\subsection{Automated Code Verification}
Complementary to semantic evaluation is the functional verification of generated code. In the context of STM32 development, this involves two layers of testing:
\begin{enumerate}
    \item \textbf{Syntactic and Static Analysis:} Static analysis examines source code without execution to identify potential vulnerabilities and coding standard violations. Tools like \texttt{flake8}, \texttt{mypy}, or \texttt{pylint} are integrated into the workflow to ensure that generated Python scripts (e.g., for dataset management or NNI optimization) are syntactically correct and follow consistent type annotations \cite{surveyOnStaticAnalysis}. This proactive layer detects errors like unused variables or null pointer risks before runtime.
    \item \textbf{Execution Loops:} The agentic framework attempts to execute generated scripts in a controlled environment. If an execution error occurs—such as a missing dependency or an invalid model layer replacement—the error log is captured and fed back to the agent as a "feedback loop." This allows the agent to self-correct the code in an iterative cycle, simulating a human developer's "trial and error" process.
\end{enumerate}

\subsection{Human-in-the-Loop Validation}
While automated loops handle many routine errors, complex architectural decisions often require human oversight. \textit{Human-in-the-loop (HITL)} patterns allow the framework to pause execution at critical junctures—such as before modifying a pre-trained model's layers or deploying a generated binary—enabling the user to approve, modify, or reject the proposed changes. This collaborative approach combines the speed of agentic automation with the safety of human expertise, ensuring that the final Edge AI solution is both optimized and trustworthy.
